\section{Spacetime and Matter}

\noindent\textbf{Exercise~\ref{ex:spacetimematter-1}.}\label{sol:spacetimematter-1}

The observer $(\del,e)$ has a frame such that
\bse
    g(e_a,e_b) = \eta_{ab}, \qand e_0(\lambda) = v_{\del,\del(\lambda)}.
\ese
We thus have
\bse
    m\big(\nabla_{v_{\gamma}}v_{\gamma}\big)^a = qF^{ac}\eta_{cb}{v_{\gamma}}^b,
\ese
and so
\bse
    \begin{split}
        m\big(\nabla_{v_{\gamma}}v_{\gamma}\big)^0 & = qF^{00}{v_{\gamma}}^0 + \sum_{\beta=1}^3 qF^{0\beta}{v_{\gamma}}^{\beta} \\
        m\big(\nabla_{v_{\gamma}}v_{\gamma}\big)^{\a} & = -qF^{\a0}{v_{\gamma}}^0 - \sum_{\beta=1}^3 qF^{\a\beta}{v_{\gamma}}^{\beta}.
    \end{split}
\ese

\bigskip
\noindent\textbf{Exercise~\ref{ex:spacetimematter-2}.}\label{sol:spacetimematter-2}

We set
\bse
    \cF^{\a} = m \big(\nabla_{v_{\del}}v_{\del}\big)^{\a} = q{F^{\a}}_b{v_{\del}}^b,
\ese
where $\cF$ is the Lorentz force, up to some factors. We can split the right-hand side into two terms, one for $b=0$ and the other for $b=\beta$. For the former we note that
\bse
    \eta^{\a\beta}E_{\beta} = \eta^{\a\beta}F_{\beta0} = {F^{\a}}_0,
\ese
so, using $\epsilon^0:(v_{\del}) = {v_{\del}}^0$, the first term is simply
\bse
    q\eta^{\a\beta}E_{\beta}\big(\epsilon^0:v_{\del}\big).
\ese
The second term need a little more work, but we start using the hint: using $g^{ab}=\eta^{ab}$ in the observers frame, we have
\bse
    (v_{\del}\times B)^{\a} = \eta^{\a\mu}\varepsilon_{\mu\rho\sig} {v_{\del}}^{\rho}B^{\sig} := \frac{1}{2}\eta^{\a\mu}\varepsilon_{\mu\rho\sig}\varepsilon^{\sig\nu\tau} {v_{\del}}^{\rho}F_{\nu\tau}.
\ese
We see that the $\eta^{\a\mu}$ tells us that $\mu=\a$ in the Levi-Civita symbol. By definition, the only non-vanishing terms on the right, then, have $\rho\neq\sig\neq\a$. Let's consider the case for $\a=1$ (the other two follow analogously)
\bse
    \begin{split}
        \frac{1}{2}\eta^{1\mu}\varepsilon_{\mu\rho\sig}\varepsilon^{\sig\nu\tau}F_{\nu\tau}{v_{\del}}^{\rho} & = \frac{1}{2}\varepsilon_{123}\big(\varepsilon^{312}F_{12} + \varepsilon^{321}F_{21}\big){v_{\del}}^2 + \frac{1}{2}\varepsilon_{132}\big(\varepsilon^{213}F_{13} + \varepsilon^{231}F_{31}\big){v_{\del}}^3 \\
        & = \frac{1}{2} \varepsilon_{123} \big( \varepsilon^{123}F_{12} + (-\varepsilon^{123})(-F_{12})\big) {v_{\del}}^2 + \frac{1}{2} (-\varepsilon_{123}) \big( (-\varepsilon^{123})F_{13} + \varepsilon^{123}(-F_{13})\big){v_{\del}}^3 \\
        & = F_{12}{v_{\del}}^2 + F_{13}{v_{\del}}^3,
    \end{split}
\ese
where we have indicated where the antisymmetries of $\varepsilon$ and $F$ have been used, and we have also used $\varepsilon_{123}=1=\varepsilon^{123}$. This generalises to
\bse
    \frac{1}{2}\eta^{\a\mu}\varepsilon_{\mu\rho\sig}\varepsilon^{\sig\nu\tau}F_{\nu\tau}{v_{\del}}^{\rho} = \eta^{\a\sig}F_{\sig\beta}{v_{\del}}^{\beta} = {F^{\a}}_{\beta}{v_{\del}}^{\beta},
\ese
where we note the fact that $F_{\a\a}=0$ due to antisymmetry.

So the second term in our Lorentz force equation is simply
\bse
    q{F^{\a}}_{\beta}{v_{\del}}^{\beta} = q\big( v_{\del} \times B \big)^{\a},
\ese
giving us
\bse
    \cF^{\a} = q\eta^{\a\beta} E_{\beta} \big( \epsilon^0:v_{\del}\big) + q\big(v_{\del}\times B\big)^{\a}.
\ese
Finally, we note that
\bse
    \frac{1}{(\epsilon^0:v_{\del})} \big(v_{\del}\times B\big)^{\a} = \big(\mathbf{v}\times B\big)^{\a},
\ese
as everywhere in the calculation above we'll get terms like
\bse
    \frac{v_{\del}^\rho}{(\epsilon^0:v_{\del})} = \frac{(\epsilon^\rho:v_{\del})}{(\epsilon^0:v_{\del})} =: \mathbf{v}^{\rho}.
\ese
This gives us the form we want on the right-hand side, i.e.
\bse
    \frac{1}{(\epsilon^0:v_{\del})}\cF^{\a} = q\eta^{\a\beta}E_{\beta} + q\big(\mathbf{v}\times B\big)^{\a}.
\ese

\bigskip
\noindent\textbf{Exercise~\ref{ex:spacetimematter-3}.}\label{sol:spacetimematter-3}

First use the antisymmetry to give
\bse
    {R^w}_{zab;c} + {R^w}_{zbc;a} + {R^w}_{zca;b} = {R^w}_{zab;c} + {R^w}_{zbc;a} - {R^w}_{zac;b}.
\ese
Now, we can use the fact that $\nabla$ is metric-compatible and the result
\bse
    g^{av}g_{vw} = \del^a_w
\ese
to give
\bse
    \begin{split}
        g^{av}g_{vw}\big({R^w}_{zab;c} + {R^w}_{zbc;a} + {R^w}_{zca;b}\big) & = {R^a}_{zab;c} + {R^a}_{zbc;a} - {R^a}_{zac;b} \\
        & = R_{zb;c} - R_{zc;b} + {R^a}_{zbc;a}.
    \end{split}
\ese
Then contract with $g^{bz}$ along with the results $R^{\,\,\,a}_{z\,\,\,bc} = -{R^a}_{zbc}$\footnote{We essentially showed this result in tutorial 9 (we showed $R_{abcd}=-R_{bacd}$).} and $R:= g^{ab}R_{ab}$ to give
\bse
    \begin{split}
        g^{bz}\big(R_{zb;c} - R_{zc;b} + {R^a}_{zbc;a}\big) & = R_{;c} - {R^b}_{c;b} - {R^{ba}}_{bc;a} \\
        & = R_{;c} - {R^b}_{c;b} - {R^a}_{c;a} \\
        & = R_{;c} - 2{R^b}_{c;b}.
    \end{split}
\ese
Finally using the fact that the Ricci tensor is symmetric,\footnote{This result is obtained from $R_{abcd}=R_{cdab}$, then setting $a=c$ and raising the first index back up.} and contracting with $g^{ac}$, we have (after relabelling)
\bse
    \big(\nabla_aG)^{ab} = {R^{ab}}_{;a} - \bigg(\frac{1}{2}g^{ab}R\bigg)_{;a} = 0.
\ese
